\uuid{M3t9}
\titre{Distance}
\auteur{bodin ; hanani ; mzari}
\organisation{exo7}
\licence{CC BY-NC-SA 4.0}
\datecreate{2020-01-20}
\niveau{L1}
\module{Géométrie}
\chapitre{Géométrie de l'espace}
\sousChapitre{Distance}
\difficulte{3}
\theme{Distance}
\exo7qcmid{lille-1.geometrie-espace.028}

\contenu{
\texte{
Pour ces questions, l'espace est muni d'un repère orthonormé direct  $(O, \vec{i}, \vec{j}, \vec{k})$.
}

\begin{question}
\qclassification{141.05}
\qtags{motcle=distance}
Soit  $P_1 : z+3=0$ et $P_2 : 2x+y+2z-1=0$ des plans et $\pi$  un plan bissecteur de $P_1$ et $P_2$, c'est-à-dire :   $M \in \pi$  si et seulement si $M$ est  à la même distance de $P_1$ et de  $P_2$. Une équation cartésienne de $\pi$ est :
\begin{answers}

       \good {$2x+y-z-10=0$ ou $2x+y+5z+8=0$}

       \bad {$x+y-z-1=0$ ou $x+y+z+1=0$}

       \bad {$2x+y+z+8=0$ ou $2x-y+5z+7=0$}

       \bad {$x+y-z-4=0$ ou $x+y+3z-8=0$}

\end{answers}
\begin{explanations}
$M(x,y,z) \in \pi \Leftrightarrow |z+3|=\frac{|2x+y+2z-1|}{3}$.
\end{explanations}
\end{question}
}
